% status: SKELETON
% Chapter 1 of the second edition. Plan: docs/STRUCTURE.md. Rules: AUTHORING.md.
\chapter{Why Inference Is Its Own Discipline}\label{ch:why-inference}

\begin{ChapterIntent}
Training produces a model once; inference runs it billions of times, one
token after another, for people who are waiting. That changes what matters:
the work is sequential rather than parallel, stateful rather than stateless,
bounded by memory rather than arithmetic, and priced per token rather than per
run. This chapter measures the difference on an ordinary machine, introduces
the four scarce resources --- memory bandwidth, memory capacity, compute and
latency --- and the two outcomes the book cares about, correctness and cost.
It then introduces the constraint ledger, the accounting every later chapter
uses to charge a technique for what it spends, and shows how a single 2026
model configuration file already names most of the book.
\end{ChapterIntent}

\OpeningSection{Sixteen conversations for the price of two}

\NapkinSection{Napkin math: training pays once, inference pays per token}

\section{What an engine does that a model does not}

\section{Four scarce resources}

\section{Two outcomes: correctness and cost}

\section{The constraint ledger}

\section{A configuration file is a table of contents}

\section{How this book is organized}

\LedgerSection

\LabSection{a fake model with real costs}

\HappenedSection{one context, four users}

\FrontierSection{}

\ExercisesSection
